\subsection{Conclusion}
\label{cic:sec:conclusion}

We present Control in Context, a mixed-methods study of how and why smart home users navigate the spectrum from physical per-device interaction to proxy-based central control. Through a survey of $N=60$ smart home owners and semi-structured follow-up interviews with $N=8$ participants, we find that smart home control is a dynamic, reflective, and often labor-intensive process rather than a static preference. Users continuously adapt their control strategies---and their device procurement decisions---as they encounter tradeoffs between convenience, reliability, precision, and cognitive load. In particular, over half of participants reported switching control schemes based on task context, and nearly all interview participants described at least one device or platform abandoned after a negative troubleshooting experience. Our work emphasizes that \textit{future systems should not only support multiple control schemes, but also actively guide users' movement across them}. Researchers should similarly reframe their evaluation metrics and system assumptions to reflect the temporal and spatial heterogeneity of real-world smart home usage. This work is an important initial step towards home management systems that are truly user-centric rather than device-centric.
